\section{空间域滤波：把邻域、噪声和边缘联系起来}
\noindent\textbf{阅读定位。}本节主要对应9月28日录音02:01:14--02:19:35。老师关注的是“为什么这样处理、得到什么、牺牲什么”，并非只背模板。以下数值题均标为\textbf{补充算例}，是为理解课堂观点编写的练习，不冒充课堂原题。公式中的图像先按实数处理，最后显示时再截断或量化。

\subsection{模板操作、相关与卷积}
设输入为 $f(x,y)$，模板为 $h(s,t)$，模板中心对应 $(s,t)=(0,0)$。对实值图像，相关操作为
\[
g_{\mathrm{corr}}(x,y)=\sum_s\sum_t h(s,t)f(x+s,y+t),
\]
而卷积为
\[
g_{\mathrm{conv}}(x,y)=\sum_s\sum_t h(s,t)f(x-s,y-t).
\]
计算卷积时，可先将模板旋转 $180^\circ$，再逐位置相乘求和。模板中心对称时，相关和卷积结果相同；非对称模板不能直接混用。方向差分模板尤其可能差一个负号。

老师在9月28日02:03:28用 $3\times3$ 邻域说明“对应元素相乘、相加，写入结果对应位置”。这个操作描述直接对应相关；课堂中常将模板滤波概括称为卷积。理解区别比纠结名称更重要：\textbf{考试给出约定，就按题目约定；自己写答案时，先说明是否翻转模板。}

\paragraph{边界为什么必须交代？}
当中心靠近图像边缘时，部分邻域落在图外。常见选择如下：
\begin{itemize}
\item 不补边，只算完整邻域：输出变小，例如 $M\times N$ 图像与 $3\times3$ 模板得到 $(M-2)\times(N-2)$ 输出
\item 零填充：容易使亮图像边缘变暗，并引入人为跳变
\item 复制边缘或镜像延拓：通常较少制造明显边缘，但镜像是否重复端点必须约定
\item 周期延拓：DFT天然对应的模型；左右或上下不连续时会出现接缝效应
\end{itemize}
老师在9月28日02:03:28提到舍弃或复制边界。实际工程中，边界占比、测量区域和后续处理会决定这种选择是否重要。计算时还应从原图读、向新图写，不能边扫描边覆盖输入，否则不再是同一个固定模板滤波。

\subsection{补充算例1：一次模板计算，为什么可能得到两个答案？}
给定邻域和非对称模板
\[
P=\begin{bmatrix}1&2&3\\4&5&6\\7&8&9\end{bmatrix},\qquad
h=\begin{bmatrix}1&0&-1\\1&0&-1\\1&0&-1\end{bmatrix}.
\]
\textbf{相关：}按位置直接相乘，得到
\[
g_{\mathrm{corr}}=(1-3)+(4-6)+(7-9)=-6.
\]
\textbf{卷积：}模板旋转后左右符号互换，得到
\[
g_{\mathrm{conv}}=(-1+3)+(-4+6)+(-7+9)=6.
\]
二者幅值相同，方向符号相反。如果任务仅取梯度幅值，这次差异不改变强度；如果要判断由暗到亮还是由亮到暗，则不能丢掉符号。

再考虑一维图像 $[10,20,30]$，模板为 $\frac13[1,1,1]$，求最左像素：
\[
\begin{aligned}
\text{零填充：}&\quad (0+10+20)/3=10,\\
\text{复制端点：}&\quad (10+10+20)/3=40/3,\\
\text{不重复端点的镜像：}&\quad (20+10+20)/3=50/3.
\end{aligned}
\]
\textbf{结论：}同一模板不保证同一数值，必须同时说明运算定义、锚点和边界处理。

\subsection{均值、高斯平滑：减噪的代价是混合相邻结构}
典型的均值模板与高斯形加权模板为
\[
h_{\mathrm{mean}}=\frac19\begin{bmatrix}1&1&1\\1&1&1\\1&1&1\end{bmatrix},\qquad
h_{\mathrm{G}}=\frac1{16}\begin{bmatrix}1&2&1\\2&4&2\\1&2&1\end{bmatrix}.
\]
右式是常用的离散高斯近似，而不是任意参数下精确采样的高斯。一般高斯权重由
\[
\widetilde h(s,t)=\exp\left[-\frac{s^2+t^2}{2\sigma_s^2}\right],\qquad
h(s,t)=\frac{\widetilde h(s,t)}{\sum_{a,b}\widetilde h(a,b)}
\]
归一化得到。这里 $\sigma_s$ 是空间尺度，不是噪声标准差。权重非负且和为1，输出是邻域灰度的加权平均，不会超出该邻域的最小值与最大值。

\paragraph{为什么强调系数和？}
设图像局部恒为 $c$，则 $g=c\sum h$。因此 $\sum h=1$ 保持常量区域亮度；$\sum h=0$ 则使常量输出为0，常见于导数算子。注意“和为1”不意味着一定平滑，锐化核也可能和为1，只是含负系数。边界处理不当时，和为1也不能保证有限图像的全局均值严格不变。课堂对应：9月28日02:05:40--02:07:55。

\paragraph{为什么能减噪？}
若 $f_i=s+n_i$，邻域内真实信号近似同一常量 $s$，噪声独立、零均值、方差均为 $\sigma_n^2$，且 $\sum_iw_i=1$，则
\[
\mathbb E[g]=s,\qquad \operatorname{Var}(g)=\sigma_n^2\sum_iw_i^2.
\]
$K$ 点均值使方差降为 $\sigma_n^2/K$，标准差降为 $\sigma_n/\sqrt K$。噪声相关时还要加协方差项；不能把这个结果无条件用于条纹噪声。高斯噪声描述\emph{噪声幅值分布}，高斯滤波描述\emph{模板权重形状}，二者不是同一个概念。

\subsection{补充算例2：比较减噪强度，再解释边缘为什么变糊}
一个平坦区域真实灰度为100，独立零均值噪声标准差为12。采用上述两个 $3\times3$ 模板。

均值滤波的输出期望仍为100，方差与标准差为
\[
\sigma_{\mathrm{mean}}^2=144\times9(1/9)^2=16,\qquad
\sigma_{\mathrm{mean}}=4.
\]
高斯形模板权重平方和为
\[
\sum_iw_i^2=\frac{4\times1^2+4\times2^2+4^2}{16^2}
=\frac{36}{256}=\frac9{64},
\]
所以
\[
\sigma_{\mathrm G}^2=144\times\frac9{64}=20.25,\qquad
\sigma_{\mathrm G}=4.5.
\]
这说明在\emph{相同的9个独立同方差观测、且真实信号恒定}的条件下，均匀平均减噪更强。但实际邻域可能跨过边界，距离中心远的点未必来自同一对象，高斯加权就提供了不同的结构保留折中。

用一维阶跃看代价：令 $f[n]=0$（$n<0$），$f[n]=100$（$n\geq0$）。三点均值输出在 $n=-2,-1,0,1$ 处依次为
\[
0,\quad \frac{100}{3},\quad \frac{200}{3},\quad100.
\]
原来一次跳变变成逐渐过渡。改用 $\frac14[1,2,1]$，相应输出为 $0,25,75,100$。中心权重较大的核，跨边界混合更轻；进一步增大平滑尺度，往往能减噪更多，却让边缘过渡更宽。

\textbf{可直接用于解释题的答案：}平滑把周围像素平均进来，平坦区的独立随机波动部分抵消；在边界处却把两个对象的灰度混合，因而降低噪声的同时损失细节。应按噪声水平和最小目标尺寸选择尺度，而非“窗口越大越好”。课堂对应：9月28日02:07:55。

\subsection{中值滤波：排除离群点，但不保证几何不变}
中值滤波把邻域排序，取中间值。奇数个元素 $a_{(1)}\leq\cdots\leq a_{(K)}$ 的中值为 $a_{((K+1)/2)}$。它对少量极端椒盐噪声有用，因为极端值被排在两端，不一定影响中间值；均值却会被极端值拉偏。

中值是非线性的，不能把它当作固定卷积核。它可以保留某些理想阶跃，但会消除小亮点、细线，改变拐角甚至移动边界。在精密尺寸测量中，不能只看“图像干净了”。课堂对应：9月28日02:17:58--02:19:35。课堂明确举了椒盐点被中值去除、但测量边界可能改变的对比情境。

\subsection{补充算例3：椒盐噪声消失了，真实细线也可能消失}
局部 $3\times3$ 邻域为
\[
P=\begin{bmatrix}100&100&100\\100&255&100\\100&100&100\end{bmatrix}.
\]
中心是一个盐噪声点。均值输出
\[
\frac{8\times100+255}{9}=\frac{1055}{9}\approx117.22,
\]
中值输出100。均值只是降低异常点强度，并把其影响传播到相邻输出；中值在这个例子中恢复了背景值。

但是，如果这个255是需要保留的真实亮点，中值同样会删除它。进一步考虑背景为0、一条宽1像素的竖直亮线为255：以亮线为中心的 $3\times3$ 窗口内只有3个255和6个0，中值为0，亮线被消除。

\textbf{非线性可用一个反例证明。}在一个三元素窗口内，取
\[
a=(0,0,100),\qquad b=(0,100,0).
\]
则 $\operatorname{med}(a)=\operatorname{med}(b)=0$，但
\[
\operatorname{med}(a+b)=\operatorname{med}(0,100,100)=100\ne0+0.
\]
因此不满足可加性。10月12日01:37:12--01:37:56老师用“中值滤波频谱特性”作为思考题：它有平滑、抑制尖点的效果，可定性说有低通倾向；\textbf{不能因此写成适用于所有输入的固定 $G=HF$ 线性传递函数。}

\subsection{导数、拉普拉斯与锐化：先分清“边缘图”和“增强图”}
一阶导数描述变化方向与变化强度。离散形式例如
\[
D_xf(x,y)=f(x+1,y)-f(x,y),\quad
D_yf(x,y)=f(x,y+1)-f(x,y),\quad
\lVert\nabla f\rVert=\sqrt{(D_xf)^2+(D_yf)^2}.
\]
梯度幅值大表示变化快；梯度方向垂直于局部等灰度线。仅求幅值会丢掉明暗变化方向。

四邻域拉普拉斯采用以下符号约定：
\[
Lf=f(x+1,y)+f(x-1,y)+f(x,y+1)+f(x,y-1)-4f(x,y),
\qquad
h_L=\begin{bmatrix}0&1&0\\1&-4&1\\0&1&0\end{bmatrix}.
\]
其系数和为0，因此恒定区域输出0；它是二阶差分，不是保留原图亮度的增强图。若用于锐化，应把带符号的导数加回原图：
\[
g=f-\lambda Lf,\qquad \lambda>0.
\]
在此负中心约定下是\emph{减}拉普拉斯；若定义的是正中心模板，符号反过来。$\lambda=1$ 时锐化核为
\[
h_{\mathrm{sharp}}=\begin{bmatrix}0&-1&0\\-1&5&-1\\0&-1&0\end{bmatrix}.
\]
八邻域版本相应是中心 $-8$、其余1的导数核，以及中心9、其余 $-1$ 的锐化核。\textbf{课堂明确比较过中心4/5与8/9的含义}：前者属于差分提边，后者含原图直流分量，用于锐化（9月28日02:11:32--02:13:21）。

\paragraph{反锐化掩蔽并不是“变得不锐”。}
先求平滑图 $\bar f=h*f$，再形成细节掩模 $m=f-\bar f$，最后
\[
g=f+km=(1+k)f-k\bar f,\qquad k\geq0.
\]
$k=0$ 不变，$k=1$ 是常见反锐化掩蔽，增大 $k$ 增强更强，也更易放大噪声、产生光晕。等效核为 $(1+k)\delta-kh$，系数和仍为1。课堂对应：9月28日02:07:55--02:10:11。\textbf{用于加回原图的掩模应保留正负号}，取绝对值适合看边缘强度，却会改变锐化结果的含义。

\paragraph{课件与教材的高提升约定不同。}
增强课件第85页写 $g_A=Af-\bar f$，其核系数和为 $A-1$，因此在 $A>2$ 时还会提高常量区域亮度。上文采用教材第183页的 $g_k=f+k(f-\bar f)$，其系数和始终为1。两者不能仅令 $A=k+1$ 就认定相等，因为平滑项系数分别为 $-1$ 和 $-k$。例如 $A=2$、$k=1$ 时才直接得到同一式子 $2f-\bar f$；如果允许整体重新缩放，则需把缩放因子也写清楚。考试应按给出的定义计算，避免只记“高提升”名称。

\subsection{补充算例4：一次锐化计算，以及一个噪声警告}
某像素中心灰度100，上下左右灰度分别为90、90、100、100，则
\[
Lf=90+90+100+100-4\times100=-20,
\qquad g=f-Lf=100-(-20)=120.
\]
如果该像素的平滑值 $\bar f=95$，采用 $k=1.5$ 的反锐化掩蔽，则
\[
m=100-95=5,\qquad g=100+1.5\times5=107.5.
\]
两个结果不同并不矛盾，因为模板与增强强度不同。若所有邻域像素都是100，拉普拉斯为0，锐化仍为100。

为什么噪声图像不宜直接二阶微分？考虑一维交替噪声 $n[j]=a(-1)^j$。一阶差分满足
\[
n[j+1]-n[j]=-2a(-1)^j,
\]
振幅变为 $2a$；二阶差分满足
\[
n[j+1]-2n[j]+n[j-1]=-4a(-1)^j,
\]
振幅变为 $4a$。这解释了快速起伏的噪声会得到很强响应。导数也响应真实边缘，\emph{不是只放大噪声而不响应信号}；困难在于它无法自动区分两者。课堂在9月28日02:15:56--02:17:58强调先估计噪声、必要时先平滑，再提边或锐化。

\paragraph{实做检查。}负值与超过255的值是正常的中间结果。若过早转无符号8位，可能截断或绕回而破坏计算。应先用浮点完成处理，再按输出用途选取截断、归一化或带符号可视化。平滑与锐化不是互相抵消的机械步骤：前者抑制不可信的快速起伏，后者增强保留下来的结构，强度必须协同选择。

\section{10月12日频域材料：小测是否纳入尚未确认}
\noindent\textbf{范围提示。}10月12日00:00:00--00:00:53宣布下周用15分钟回答一两个小问题，覆盖“前面讲过”的内容，强调理解、不要求背记。原话未明确界定是否包含当天随后讲授的频域材料。因此本节是\textbf{单独列出的频域复习与扩展}，不能据此宣称它一定在本次小测范围内。

\subsection{空间频率：快慢是灰度随位置变化的快慢}
低频对应缓慢变化、大尺度结构；高频对应快速变化、尖锐边界和细纹理。它不是“灰度越高、频率越高”，也不是“所有噪声都是高频”。均匀的白图只有直流分量；周期性干扰可能集中在中频；真实细节也可能是高频。

对 $M\times N$ 图像，采用前向不归一化、逆向除以 $MN$ 的约定：
\[
F(u,v)=\sum_{x=0}^{M-1}\sum_{y=0}^{N-1} f(x,y)
\exp\left[-\mathrm i2\pi\left(\frac{ux}{M}+\frac{vy}{N}\right)\right],
\]
\[
f(x,y)=\frac1{MN}\sum_{u=0}^{M-1}\sum_{v=0}^{N-1} F(u,v)
\exp\left[\mathrm i2\pi\left(\frac{ux}{M}+\frac{vy}{N}\right)\right].
\]
这里 $\mathrm i^2=-1$。直流项 $F(0,0)=\sum f$，所以平均灰度是 $F(0,0)/(MN)$。幅度谱为 $|F|$，相位为 $\arg F$；一般需要完整复数频谱才可无损重建，只有显示出来的幅度图不够。

二维DFT可先逐行、再逐列计算。DFT是变换，FFT是计算DFT的快速算法；DCT是另一种离散余弦变换，不是DFT的别名。二维FFT典型复杂度为 $O(MN(\log M+\log N))$，不是指数复杂度。课堂对应：10月12日00:22:06--00:26:30；这些区分用于避免把口语或转写中的混称记成定义。

\paragraph{频谱图怎样读？}
默认DFT的零频在数组原点；居中重排后，零频在中间，越远通常表示越高的频率。居中重排只改变存储位置，既不是滤波，也不是必要的数学变换。滤波器必须与当前频谱排列一致，居中过的频谱在逆变换前要逆向重排。显示频谱常用 $\log(1+|F|)$ 压缩动态范围，但滤波应乘在\emph{复数频谱}上，而不是乘在这张对数显示图上。课堂对应：10月12日00:26:58--00:28:17。

\subsection{补充算例5：四点DFT看清直流、条纹与高通}
把一维序列视为图像的一行，令 $f=[2,0,2,0]$，长度 $N=4$。直接计算：
\[
\begin{aligned}
F[0]&=2+0+2+0=4,\\
F[1]&=2+2e^{-\mathrm i\pi}=0,\\
F[2]&=2+2e^{-\mathrm i2\pi}=4,\\
F[3]&=2+2e^{-\mathrm i3\pi}=0.
\end{aligned}
\]
因此 $F=[4,0,4,0]$。它由均值1和交替分量 $[1,-1,1,-1]$ 组成。仅保留直流，反变换得到 $[1,1,1,1]$；去掉直流、保留 $F[2]$，得到 $[1,-1,1,-1]$。

\textbf{解释：}低通保留背景，抹去快速变化；高通得到零均值的有符号细节，并不自动得到适于直接显示的普通灰度图。把高通结果直接截到 $[0,255]$ 会丢掉负细节。如果想锐化，应按需要把细节加回原图。

验证能量：空间域 $\sum|f[n]|^2=8$，频域在本归一化约定下为
\[
\frac14\sum_k|F[k]|^2=\frac14(16+16)=8.
\]
一般有 $\sum|f|^2=\frac1{MN}\sum|F|^2$。能量比例应以\textbf{幅值平方和}计算。老师在10月12日00:52:42--00:54:09展示某图低频集中并报出98.8\%，但记录描述的是幅度和统计；不能将其推广为所有图像的能量比例，也不能把幅度和与严格能量定义混同。

\subsection{卷积定理：同一线性处理的两种表达}
\[
g=f*h\quad\Longleftrightarrow\quad G(u,v)=F(u,v)H(u,v).
\]
$h$ 是空间域冲激响应，$H$ 是它的频域传递函数。流程为：输入与边界准备，DFT，乘 $H$，逆DFT，取有效区域。课堂对应：10月12日00:17:56--00:19:48。

\textbf{注意周期性。}同尺寸DFT相乘后逆变换对应循环卷积。若要实现 $M\times N$ 图像与 $P\times Q$ 核的完整线性卷积，至少补零到 $(M+P-1)\times(N+Q-1)$，同时处理好核原点位置，最后按需要裁剪。若忽略补边，图像一侧的内容可能绕到另一侧。

\subsection{低通、高通、截止参数与振铃}
以居中频谱坐标为例，令 $D(u,v)$ 为到零频中心的距离，$D_0>0$ 为尺度参数。常用低通形式为
\[
H_{\mathrm I}(D)=\begin{cases}1,&D\leq D_0,\\0,&D>D_0,\end{cases}
\quad
H_{\mathrm B}(D)=\frac{1}{1+(D/D_0)^{2n}},\ n\geq1,
\quad
H_{\mathrm G}(D)=e^{-D^2/(2D_0^2)}.
\]
这三个式子为常用数字图像处理约定，不同教材参数定义可能不同。特别地，巴特沃斯在 $D_0$ 处为 $1/2$，这里的高斯在 $D_0$ 处为 $e^{-1/2}$，并不具有相同的“半幅截止点”。比较效果应注意实际带宽，而不仅比较同名参数。

对应高通可定义为 $H_{\mathrm{HP}}=1-H_{\mathrm{LP}}$。低通中 $D_0$ 越小，通常模糊越强；高通中 $D_0$ 越大，低频被去掉得越多。纯高通令直流通过率为0，因此输出整体平均值为0，不能与保亮度锐化混为一谈。课堂对应：10月12日00:39:56--00:44:42、01:04:05--01:07:27。

\paragraph{振铃不是普通的模糊。}
理想滤波器在频域突然截断，对应空间响应有振荡旁瓣；与强边缘卷积后会产生明暗交替的波纹、过冲或欠冲，这称为振铃。连续一维理想低通对应sinc型响应，连续二维圆形理想低通对应径向Bessel型响应，不宜把所有维度都说成同一个一维函数。课堂明确展示理想低通产生重影，并在10月12日00:47:31--00:51:49称之为振铃。

巴特沃斯过渡较平缓，但阶数 $n$ 增大时越来越接近理想截止，仍可能振铃；不能背成“巴特沃斯绝无振铃”。理想高通也可能振铃。理想连续高斯没有这种振荡旁瓣，但有限离散实现仍须注意截断与边界伪影。选择平滑滤波器时要同时权衡噪声、细节和伪影。课堂对应：10月12日01:13:42--01:15:53。

\subsection{补充算例6：同样的参数，不同的滤波器通过多少？}
设 $D_0=10$，巴特沃斯阶数 $n=2$。在距离零频 $D=20$ 的位置，原频谱系数为 $F=8+6\mathrm i$，其幅值为10。
\[
\begin{aligned}
H_{\mathrm I}(20)&=0,\qquad &|G_{\mathrm I}|&=0,\\
H_{\mathrm B}(20)&=\frac1{1+2^4}=\frac1{17},\qquad
& G_{\mathrm B}&=\frac8{17}+\frac6{17}\mathrm i,\\
H_{\mathrm G}(20)&=e^{-2}\approx0.1353,\qquad
& |G_{\mathrm G}|&=10e^{-2}\approx1.3534.
\end{aligned}
\]
巴特沃斯输出幅值为 $10/17\approx0.5882$。滤波乘数为非负实数时，非零系数的相位保持不变，幅值按比例缩小。对应巴特沃斯高通在此位置的通过率为 $16/17$，而不是 $1/17$。

在 $D=10$ 处，巴特沃斯通过率为 $1/2$，高斯为约0.6065；由此可见同一 $D_0$ 不意味着完全相同的过渡带。若把巴特沃斯阶数提高到4，在 $D=20$ 处通过率为 $1/(1+2^8)=1/257$，截止更陡，但应警惕振铃。

\textbf{课堂实例与练习的边界：}老师在10月12日00:41:22明确讲过“半径30的圆形低通”造成图像变糊。上面的 $D_0=10$ 和复数系数是补充计算，并非该课堂图像的真实数据。

\subsection{周期噪声：窄陷波通常比整圈带阻更有选择性}
若噪声近似为
\[
n(x,y)=A\cos\left[2\pi\left(\frac{u_0x}{M}+\frac{v_0y}{N}\right)+\phi\right],
\]
则在频域中主要出现于一对相反频率 $(u_0,v_0)$ 与 $(-u_0,-v_0)$。实值图像满足共轭对称性，滤除时通常成对处理，以保持反变换结果为实数。

\begin{itemize}
\item 带通：保留某一频带，可突出特定尺度纹理
\item 带阻：抑制某一频带，径向带阻会删除一个环带内的多种方向信息
\item 陷波：只抑制若干局部频率区域，适合位置明确的干扰峰；通常比整圈带阻损伤更小
\end{itemize}
但频谱亮点可能是真实周期纹理，不能见亮就删。条纹不严格周期、有限图像截断和频率不对齐会使能量泄漏成一小团，需要合适宽度、平滑过渡，而不一定只清零一个像素。课堂明确案例：10月12日01:08:15--01:11:18，通过分析条纹图像中的对称亮点，抑制周期干扰。

\subsection{补充算例7：准确找出条纹频率，而不是整幅图都做低通}
设 $64\times64$ 图像，$x$ 表示行坐标、$y$ 表示列坐标，观测为
\[
f(x,y)=100+20\cos(2\pi y/8).
\]
条纹沿列坐标每8像素重复，因而在图像中表现为竖条纹。每行包含 $64/8=8$ 个周期。未居中的DFT有
\[
F(0,0)=100\times4096=409600,
\]
以及共轭频率点
\[
(0,8),\quad(0,56),\qquad
F(0,8)=F(0,56)=\frac{20}{2}\times4096=40960.
\]
居中后它们位于中心左右各8个频率格的位置。在这个理想模型中，把这一对系数设为0，保留直流，逆变换恰好得到灰度100的常量图。

为什么不直接用半径小于8的低通？这个合成例只有背景，因此也能奏效；真实图像还含许多合法边缘和纹理，低通会同时删除全部较高频率。窄陷波只针对干扰附近，更有机会保住其他方向与频率的内容。若真实纹理与噪声恰好同频，仅凭该频谱不能无损分离二者，需要额外先验或采集端改善。

\subsection{同态滤波：把乘性的照明问题变成可处理的加法问题}
常用模型为 $f(x,y)=i(x,y)r(x,y)$，其中 $i$ 表示照明，$r$ 表示反射特性。它是一个模型假设；典型假设是照明缓变、反射细节变化较快，但两者频谱并非严格不重叠。

对正值图像取自然对数：
\[
z=\ln f=\ln i+\ln r,\qquad
Z=\mathcal F\{z\},\quad S=HZ,\quad
s=\mathcal F^{-1}\{S\},\quad g=e^s.
\]
零值输入可先用 $\max(f,\varepsilon)$ 保护。若使用 $\ln(1+f)$ 则相应逆操作应为 $e^s-1$，但它不再是对乘性模型的严格加法分解，不要在推导中悄悄替换。

常见平滑传递函数为
\[
H(D)=(\gamma_H-\gamma_L)
\left[1-\exp\left(-c\frac{D^2}{D_0^2}\right)\right]+\gamma_L,
\quad c>0,\quad 0<\gamma_L<1<\gamma_H.
\]
低频趋近 $\gamma_L$，适度压缩照明起伏；高频趋近 $\gamma_H$，增强细节。它不是把低频全部删除的纯高通。课堂对应：10月12日01:17:43--01:24:33。老师在01:25:13举室外拍室内设备、管路细节改善的例子；这说明增强已有信息，不意味着能恢复饱和或完全未记录的细节。

\subsection{补充算例8：用亮度比检验同态处理的目标}
先作一个\emph{理想分离近似}：照明对数完全在低频，反射对数完全在高频。取 $\gamma_L=0.5$、$\gamma_H=1.5$，则
\[
g\approx\exp(0.5\ln i+1.5\ln r)=i^{0.5}r^{1.5}.
\]
两位置反射率相同、照明分别为4和16，原图亮度比为 $16/4=4$；处理后比为
\[
\frac{16^{0.5}r^{1.5}}{4^{0.5}r^{1.5}}=2.
\]
照明差异被压缩。另有两位置照明相同，反射率分别为0.25与0.5，原比值为2，处理后为
\[
\frac{i^{0.5}(0.5)^{1.5}}{i^{0.5}(0.25)^{1.5}}=2^{1.5}\approx2.828.
\]
反射对比得到增强。这个计算解释了“压低频、提高频”，不是实际图像逐像素必然遵守的精确分解，因为真实照明和反射频谱会重叠。

若使用上述平滑函数，取 $c=1$、$D=D_0$，则
\[
H(D_0)=(1.5-0.5)(1-e^{-1})+0.5\approx1.1321.
\]
故过渡频率既不是简单乘0.5，也不是简单乘1.5。参数越激进，越要检查暗部噪声、光晕与颜色失真。最终显示还可能需要动态范围映射；这一步不应被误认为模型本身的一部分。

\subsection{把方法选择写成有理由的答案}
面对“图像不好，怎么处理”的问题，建议按以下顺序组织，而不是罗列所有算法：
\begin{enumerate}
\item \textbf{描述缺陷和证据：}是随机细颗粒、极端椒盐点、周期条纹、照明不均，还是只是对比不足？看空间分布、局部统计，必要时看频谱
\item \textbf{说明假设：}例如局部真实图像平滑、噪声近似独立，或者条纹频率与目标结构可区分。若假设不成立，方法为何失效
\item \textbf{选择处理：}平坦区随机噪声可试均值或高斯；稀疏脉冲可试中值；定位明确的周期干扰可试成对陷波；乘性照明起伏可试同态；需要更清晰时在减噪后适度锐化
\item \textbf{说明代价与验证：}检查细线、边缘位置、过冲、振铃和计算时间。精密测量要验证几何误差，视觉展示要检查可见伪影；不能仅凭“看起来更锐”判断成功
\end{enumerate}
老师在10月12日01:40:54--01:46:15用机器人实时处理举例，提醒先查噪声源和算力，而不是只在PC上证明算法可运行；在01:55:48--01:58:23提醒不要对强噪声图直接上拉普拉斯。\textbf{更亮的照明、更合适的曝光或更稳定的采集，有时比后期堆叠滤波更有效；具体原因仍应通过现场观测验证。}

